\documentclass[10pt,twocolumn]{article}

\usepackage[margin=0.55in,columnsep=0.22in]{geometry}
\usepackage[T1]{fontenc}
\usepackage[utf8]{inputenc}

\setlength{\parindent}{0pt}
\setlength{\parskip}{2pt}
\let\olditemize\itemize
\let\oldenditemize\enditemize
\renewenvironment{itemize}{\olditemize
  \setlength{\itemsep}{1pt}
  \setlength{\parskip}{0pt}
  \setlength{\parsep}{0pt}
  \setlength{\topsep}{2pt}
}{\oldenditemize}

\begin{document}
\footnotesize
\twocolumn[
\begin{center}
{\Large \textbf{Thesis Talk Outline}}\\
{\normalsize Evaluating Language Models for Coding, Mathematics, and Reasoning}
\end{center}
\vspace{0.5em}
]

\section*{Core Narrative}

\begin{itemize}
  \item Main message: evaluations are not just leaderboards. They are lenses for understanding model behavior, and the lenses we choose shape what the field builds next.
  \item Start from the post-GPT-4 moment: models looked broadly capable, familiar static benchmarks were saturating, and the strongest systems were increasingly closed.
  \item Explain why evaluation had to become more diagnostic: a useful benchmark should reveal meaningful shortcomings, explain model behavior, and suggest what to build next.
  \item Thesis arc: choose faithful metrics, analyze the behavior behind the scores, then use those insights to guide future AI software engineering systems.
\end{itemize}

\section*{Chapter 1: Introduction}

\begin{itemize}
  \item Open with the historical shift: GPT-3 and Codex were impressive but brittle; ChatGPT and GPT-4 changed the community's expectations.
  \item Emphasize that the surprising part was not only better scores, but broader apparent capability: coding games, proofs, reasoning demos, and general assistant behavior.
  \item Introduce the central question: what should evaluation do when frontier systems improve quickly and are difficult to understand from the outside?
  \item Preview the three lessons:
  \begin{itemize}
    \item Lesson 1: focusing on the right metric without bias is crucial.
    \item Lesson 2: looking beyond the score reveals hidden nuances.
    \item Lesson 3: what we measure shapes what we build.
  \end{itemize}
  \item Transition: Chapters 2--4 each develop one lesson through coding, math, and reasoning case studies.
\end{itemize}

\section*{Chapter 2: Right Metrics Without Bias}

\textbf{Chapter setup.}
\begin{itemize}
  \item Frame as the ``what should we measure?'' chapter.
  \item Split the chapter into two kinds of work: new metrics for missing capabilities, and controls for confounds that distort existing metrics.
\end{itemize}

\textbf{CRUXEval: introduce the benchmark.}
\begin{itemize}
  \item Motivation: code generation benchmarks such as HumanEval and MBPP tested writing code from specs, but programming also requires reading code, understanding it, and simulating execution.
  \item Introduce CRUXEval as a benchmark for Code Reasoning, Understanding, and eXecution Evaluation.
  \item Explain the two tasks:
  \begin{itemize}
    \item Output prediction: given a short Python function and input, predict the output.
    \item Input prediction: given a short Python function and desired output, find an input that produces it.
  \end{itemize}
  \item Talking point: the examples are intentionally simple for humans, because the question is whether models have a basic mental model of execution at all.
\end{itemize}

\textbf{CRUXEval: benchmark construction.}
\begin{itemize}
  \item Generate candidate functions and inputs with Code Llama 34B using Python standard-library functions as seeds.
  \item Execute functions to obtain deterministic input-output pairs.
  \item Filter aggressively: short programs, no syntax errors, bounded runtime, low arithmetic burden, deterministic behavior, no imports or side effects.
  \item Select 800 examples after checking that the dataset was large enough to preserve important model comparisons under bootstrap sampling.
  \item Speaking emphasis: the generate-and-filter recipe made it possible to find simple examples that frontier models still failed.
\end{itemize}

\textbf{CRUXEval: evaluation.}
\begin{itemize}
  \item Evaluate open and closed coding models, including StarCoder, Mistral, WizardCoder, Phi, Phind, Code Llama, DeepSeek Coder, GPT-3.5, and GPT-4.
  \item Use execution-based correctness for both tasks: generated output must match actual output; generated input must satisfy the assertion.
  \item Report pass@1 and pass@5 with bootstrap intervals and paired comparisons.
  \item Key result: GPT-4 was strongest, but still failed many simple examples; the best open models failed over half the time.
  \item Modern framing: today's models largely saturate the original benchmark, which reinforces that the benchmark named a capability at the right moment.
\end{itemize}

\textbf{CRUXEval: impact of anonymizing functions.}
\begin{itemize}
  \item Motivation: rule out the concern that models were exploiting memorized or semantically suggestive variable names.
  \item Ablation: replace variable names with anonymous identifiers such as \texttt{x1, x2, ...} while preserving function semantics.
  \item Evaluation: test Code Llama 7B, 13B, and 34B with the same prompts on the anonymized subset.
  \item Takeaway: relative performance stayed similar, strengthening the case that the benchmark measured execution reasoning rather than memorization or name cues.
\end{itemize}

\textbf{SlopCodeBench.}
\begin{itemize}
  \item Introduce the motivation: as models moved into longer software workflows, correctness alone no longer captured whether generated code remained maintainable.
  \item Benchmark setup: study iterative coding trajectories rather than one-shot function generation.
  \item Metrics: verbosity and structural erosion as proxies for code slop.
  \item Evaluation: track correctness, cost, and quality over time, then calibrate against human repository histories.
  \item Takeaway: models can keep making apparent task progress while degrading code structure, which matters for real software engineering.
\end{itemize}

\textbf{IneqMath.}
\begin{itemize}
  \item Introduce the task: solving inequality proof problems.
  \item Dataset/metric setup: compare final-answer accuracy with step-wise soundness of the reasoning trace.
  \item Key result: a model can frequently land on the correct final answer while producing invalid proof steps.
  \item Takeaway: grading only the final answer can support the wrong conclusion about mathematical reasoning ability.
\end{itemize}

\textbf{Counterfeit Conundrum.}
\begin{itemize}
  \item Introduce the problem: evaluating whether code models understand incorrect code.
  \item Benchmark idea: create or identify ``counterfeit'' programs that are wrong but plausible, avoiding trivial negative examples.
  \item Evaluation target: correctness checking, execution, and program repair can all look easier when negatives are too obvious.
  \item Takeaway: negative sample construction is a confound; filtering trivial incorrect samples changes how we interpret model accuracy.
\end{itemize}

\textbf{LiveCodeBench.}
\begin{itemize}
  \item Introduce the problem: static public coding benchmarks are vulnerable to contamination and benchmark overfitting.
  \item Benchmark construction: collect time-indexed coding tasks and evaluate models according to problem release dates.
  \item Evaluation: compare performance before and after model data cutoffs; include a broader suite of code-related tasks.
  \item Takeaway: accounting for time and contamination can substantially change model comparisons.
\end{itemize}

\textbf{Chapter transition.}
\begin{itemize}
  \item Chapter 2 changes the question from ``what score did the model get?'' to ``what did we measure, and what biases might distort it?''
  \item Chapter 3 asks what model behavior produced the score.
\end{itemize}

\section*{Chapter 3: Beyond the Score}

\textbf{Chapter setup.}
\begin{itemize}
  \item Frame as the diagnosis chapter. Aggregate scores rank systems, but deeper analysis reveals which capabilities transfer, which failures are distinct, and where future methods should intervene.
\end{itemize}

\textbf{CRUXEval quantitative analysis.}
\begin{itemize}
  \item Compare CRUXEval against HumanEval-style code generation.
  \item Key observation: distilled models improved strongly on HumanEval but not similarly on CRUXEval.
  \item Interpretation: fine-tuning on natural-language-to-code data can improve code generation without improving underlying execution reasoning.
  \item Mention the correlation between input and output prediction, and contrast it with the weaker relationship to HumanEval for distilled models.
\end{itemize}

\textbf{CRUXEval prompting and fine-tuning analysis.}
\begin{itemize}
  \item Chain-of-thought: helped GPT-4 more than other models; helped output prediction more reliably than input prediction.
  \item Fine-tuning: input-output assertion data improved scores, but did not make the task solved.
  \item Qualitative failures: even GPT-4 made mistakes on short programs, sometimes due to brittle reasoning over variables, order, or execution state.
  \item Takeaway: the score exposed a gap, but the analyses explained what kinds of code reasoning remained hard.
\end{itemize}

\textbf{Counterfeit Conundrum analysis.}
\begin{itemize}
  \item Ask what kinds of incorrect programs confuse models.
  \item Analyze whether counterfeits depend on problem difficulty or generator strength.
  \item Takeaway: subtle incorrectness is not merely a weak-model or hard-problem artifact; it is a general phenomenon in code understanding.
\end{itemize}

\textbf{LiveCodeBench analysis.}
\begin{itemize}
  \item Use the benchmark to examine HumanEval overfitting, post-training effects, and open versus closed model gaps.
  \item Speaking emphasis: high scores on familiar benchmarks can be ambiguous without time-controlled analysis.
  \item Takeaway: deeper analysis can distinguish broad capability from adaptation to specific benchmark distributions.
\end{itemize}

\textbf{LINC.}
\begin{itemize}
  \item Introduce the task: natural-language logical reasoning over premises and conclusions.
  \item Method: translate natural language into first-order logic, then run a symbolic prover.
  \item Evaluation: compare LINC against chain-of-thought across logical reasoning benchmarks.
  \item Surface result: LINC often outperformed chain-of-thought.
  \item Deeper result: LINC and chain-of-thought had distinct failure modes; LINC could fail at formalization while chain-of-thought failed at natural-language deduction.
  \item Takeaway: method comparisons should examine failure structure, not just final accuracy.
\end{itemize}

\textbf{Chapter transition.}
\begin{itemize}
  \item Chapter 3 turns benchmark results into scientific understanding.
  \item Chapter 4 asks what happens when these capabilities and failure modes become shared targets for the community.
\end{itemize}

\section*{Chapter 4: What We Measure Shapes What We Build}

\textbf{Chapter setup.}
\begin{itemize}
  \item Frame as the impact chapter. A benchmark need not remain difficult forever; its lasting value can be naming a capability and making it actionable.
\end{itemize}

\textbf{CRUXEval impact.}
\begin{itemize}
  \item Introduce the claim: CRUXEval made execution reasoning visible as a capability beyond generation.
  \item Benchmark follow-ups: extensions to more languages, runtime state prediction, semantic reasoning, and other code-reasoning tasks.
  \item Training follow-ups: works used execution traces, input-output prediction, and trace rationales as supervision.
  \item Model reporting: CRUXEval became a diagnostic benchmark in later code model releases.
  \item Takeaway: once execution reasoning was measurable, it became something model builders could train for and report on.
\end{itemize}

\textbf{LiveCodeBench impact.}
\begin{itemize}
  \item Introduce the methodological contribution: time-indexed, contamination-aware evaluation.
  \item Explain why it mattered: coding benchmarks need maintenance, not just one fixed release.
  \item Takeaway: LiveCodeBench changed how model comparisons should be interpreted in a fast-moving field.
\end{itemize}

\textbf{LeanDojo.}
\begin{itemize}
  \item Introduce the setting: theorem proving with language models in Lean.
  \item Benchmark/toolkit construction: extract data from Lean, expose premise context, and allow interaction with the Lean environment.
  \item ReProver: retrieval-augmented tactic generation via premise selection and context augmentation.
  \item Impact: LeanDojo became infrastructure that enabled future theorem-proving research.
  \item Takeaway: an evaluation can shape a field by becoming the environment researchers build on.
\end{itemize}

\textbf{ProofOptimizer.}
\begin{itemize}
  \item Introduce the problem: proofs should be correct, but also shorter and more maintainable.
  \item Method: use the Lean compiler as a verifier while models propose proof transformations.
  \item Training/evaluation: proof length becomes a measurable optimization target.
  \item Takeaway: changing the metric changes what systems learn to improve.
\end{itemize}

\textbf{Chapter transition.}
\begin{itemize}
  \item Synthesize the three lessons: choose faithful metrics, analyze behavior behind scores, and recognize that evaluations shape future systems.
  \item Chapter 5 uses these lessons to articulate what AI software engineering should measure and build next.
\end{itemize}

\section*{Chapter 5: Future Research and Remaining Challenges}

\textbf{Chapter setup.}
\begin{itemize}
  \item Frame Chapter 5 as a synthesis rather than a separate survey.
  \item The chapter asks what the earlier evaluations imply about the future of AI for software engineering.
\end{itemize}

\textbf{Tasks in AI software engineering.}
\begin{itemize}
  \item Introduce the taxonomy as a guardrail against overgeneralizing from one benchmark.
  \item Measures: scope, logical complexity, and level of human intervention.
  \item Example: CRUXEval isolates execution reasoning at function-level scope, but says little about project-level codebase reasoning.
  \item Use the taxonomy to broaden the space beyond code generation: transformation, testing, analysis, maintenance, scaffolding, and formal verification.
\end{itemize}

\textbf{Challenges.}
\begin{itemize}
  \item Evaluation and benchmarks: continuing Lesson 1, future benchmarks need task diversity, contamination controls, and construct validity.
  \item Effective tool usage: LeanDojo, ProofOptimizer, and SlopCodeBench all point toward systems that must use feedback, tests, compilers, and tools.
  \item Long-horizon code planning: SlopCodeBench and ProofOptimizer show that locally correct artifacts can still be poor, verbose, or hard to maintain.
  \item Semantic understanding: CRUXEval and Counterfeit reveal function-level weaknesses in subtle code understanding; future systems need this at codebase scale.
  \item High logical complexity and OOD domains: as simple benchmarks saturate, evaluations should continue moving toward capabilities that remain hard.
\end{itemize}

\textbf{Paths forward.}
\begin{itemize}
  \item Data: collect richer code data, including traces, types, call graphs, coverage, and invariants. Connect this naturally to CRUXEval follow-up work on execution traces.
  \item Training: build executable environments for code RL, echoing the role of Lean feedback in ProofOptimizer but extending it to broader software engineering.
  \item Adaptation: help models learn from specialized and changing codebases rather than assuming static training data covers every API and convention.
  \item Inference: use retrieval, code navigation, SWE tools, and neurosymbolic methods to ground model reasoning in executable evidence.
  \item Synthesis: the next stage of AI for code should be feedback-rich, tool-aware, and evaluated at the level of real software workflows.
\end{itemize}

\section*{Conclusion}

\begin{itemize}
  \item Return to the central thesis: a good evaluation reveals a meaningful capability or failure mode, teaches us why it happens, and helps the field decide what to build next.
  \item Final recap: metrics determine what we can see; deeper analysis determines what we can understand; benchmarks and tools determine what future systems learn to optimize.
  \item Closing line: as models keep improving, the most useful evaluations will be those that remain scientifically diagnostic rather than merely difficult.
\end{itemize}

\end{document}
